\section{One length for the spectral window and the ball}
\label{sec:h-cross}

The stitch spacing \eqref{eq:ell} is the only transverse scale in the
chain construction.
A random walk with in-plane diffusion constant \(D_\parallel\) and
transverse constant \(D_\perp\) in the sheet-index coordinate has return
probability
\begin{equation}
\label{eq:Zt}
Z(t)=(4\pi D_\parallel t)^{-1}(4\pi D_\perp t)^{-1/2}
\end{equation}
throughout
\(\ell_W^2/D_\parallel\ll t\ll (S\ell_W)^2/D_\perp\).
The spectral dimension \(-2\,d\ln Z/d\ln t\) equals \(3\) inside that
window, equals \(2\) at shorter time, and returns to \(2\) after the ring
is explored.
The constant \(D_\perp\) absorbs one factor of \(\ell_W^2\) if the walk is
written in a coordinate where each stitch is a unit step.

Graph distance is a different functional.
A transverse coordinate step of size \(1\) costs graph length \(\ell_W\).
The radius where \(d_{\mathrm{eff}}(r)=d\ln V/d\ln r\) crosses \(2.5\)
therefore tracks the same \(\ell_W(\lambda)\) that sets the window of
\eqref{eq:Zt}, and the ratio of the principal radii of the ball tracks
\(\ell_W\) as well (Proposition~\ref{prop:metric}).
You could regard a ball whose crossing radius follows \(\ell_W\) while its
axis ratio does not as evidence that the volume exponent has a source
other than this homogenization.
